% !TeX root = ../papers/nc-p3.tex
\section{良い還元による対称性を仮定しない閉点の全次数認証}
\label{sec:unrestricted-good-reduction}
\begin{proposition}\label{prop:unrestricted-good-reduction}
離散付値環 $V$ 上の四つの自由な階数4の空間のテンソル $W$ を取る。
分数体上の二つの隣接切断が階数6で、幾何学的剰余体上の窓が
標準14641型quadratic Koszul左右AS正則なら、幾何学的分数体上でも同じ性質が成立する。
各始点、各次数の元の経路も二次双対の経路も $V$ 上自由である。
特に剰余標数が奇数で、剰余窓に三色ずつのWalsh表示があり符号平方行列が階数2なら適用できる。
\end{proposition}
\begin{proof}
特殊点で各隣接切断の階数は6である。従って単元な六次小行列があり、
分数体で七次小行列が全て零なので、関係像は直和因子で剰余と可換である。
四位置の直和を取った二次双対を $B$ とする。各経路加群は有限表示である。
三次にはテンソルの一文字縮約から双対次元の下界4があり、四次にはテンソル自身から
下界1がある。実際、一文字切断の非零四次小行列は特殊点から持ち上がり、
各縮約は左右の関係挿入の交叉に属する。四次でも三位置の交叉に $W$ が入る。
特殊点で両下界が等号なので分数体でも等号である。五次は剰余で零ゆえ中山の補題で零であり、
以後も生成次数から零である。各次数で剰余次元と分数体次元が等しいことから
$B$ は次数ごとに自由で、最高次への左右の完全対合は単元行列式により持ち上がる。

半単純な頂点環の相対版として、$V^4$ に関する $B$ の被約bar複体を取る。
内部次数 $n$ の項は正次数の経路加群のテンソル積の有限直和であり、有限自由である。
剰余でKoszul性によりホモロジー次数 $i\ne n$ で完全である。
各微分で剰余階数の小行列式を分数体へ移すと、隣り合う階数の和は項の次元以上、
複体条件から以下であり、完全性が従う。この議論を各 $n$ に適用するので、
一様な有限検査次数を必要としない。従って幾何学的分数体上の $B$ とその二次双対はKoszulで、
Frobenius対合から左右AS正則、逆Hilbert行列から標準成長となる。
元の経路加群でも特殊・一般の次元が一致するため自由である。

最後に奇標数での入力を確認する。定理\ref{thm:walsh-presentation-space}の
九つの相異なる超平面の付値による平方類独立性と有限平坦な平方写像は奇標数でそのまま成立する。
三色ずつの双対の完全対合は定理\ref{thm:walsh-dual-frobenius}の
非退化対角形式と連結な台の証明で全点に成立する。
この既約表示族は定理\ref{thm:rank-twelve-quantum-limit}の量子点を含む。
そこに用いた基底変換は $2$ の冪と四乗根の単元だけで行え、六つの量子関係の順序単項式は
奇標数でも基底となる。従って同じ内部次数別bar複体の階数比較で、
この標数の幾何学的一般点は正則である。各次数の元の経路次元は任意の特殊点で標準値以上となる。

定理\ref{thm:walsh-normal-quadrics}の階数2での境界間移送は全等方平面の交叉次元から定まり、
個別の非零成分による除算を要しない。直交移送の行列式の符号も奇標数で保たれる。
実際、向きを検算した量子点の係数 $(0,1,i),(i,0,i),(-1,1,0)$ を
$\mathbb Z[i,1/2]$ 上に取る。そこで符号平方行列の二次小行列に単元があり、
どの剰余でも左右の境界写像が同じ核を持つ全射である。
従って移送は各奇素点で正則に定まり、分数体での行列式 $-1$ に特殊化は可換である。
各幾何学的奇標数の階数2開集合は既約で、この符号が一定となる。
分裂四元数の乗法、$\mathrm{PGL}_2\simeq\mathrm{SO}_3$、四正規二次式の交換則も
$2$ が可逆なら成立する。その中心平方による同定理の上界は一般点からの下界と一致する。
中心多項式環上の自由性と左右線型複体の証明により全次数の正則性を得る。
\end{proof}

\begin{theorem}\label{thm:regular-53-adic-ball}
定理\ref{thm:walsh-full-component}の標数53のテンソルを $w$ とする。
二隣接階数6のアフィンスキーム $Z$ で $w$ に還元される
$\mathbb Z_{53}$ 点は、76個の自由なテンソル係数の剰余類の直積で表示される。
その全ての分数体窓は標準Koszul左右AS正則である。この球には
$\overline{\mathbb Q}$ 上で定義される正則窓が十五次元成分の中でZariski稠密にある。
最小周期4、反対階数16、自明な射影安定化群を同時に課しても稠密である。
\end{theorem}
\begin{proof}
命題\ref{prop:twenty-conormal}の二十個の線型恒等式は
六次単元小行列のSchur分解と有限エタールな交点の番号付けによる。
横断的な六双一次式の全交点は剰余で二十単純点なので、射影性とJacobi判定で
第一切断の階数6軌跡の相対近傍でも有限エタールである。
従って同命題の消去は $\mathbb Z_{(53)}$ 上でも成立する。
剰余で独立な180個の微分を選び、対応する180係数について単元小行列を取る。
これらの方程式の零点は相対的に滑らかで、選ばなかった76係数への射影はエタールである。
実際、全方程式は局所エタールに180個で生成され、独立な180式は余接空間で
この生成系と可逆な係数行列で結ばれる。従って局所環で同じイデアルを生成する。
Henselの補題で各自由係数の指定にただ一つの持上げがある。

具体的に剰余窓の $m$ の三行は $(4,5,16),(23,-3,19),(-21,4,-1)$ である。
\path{computations/verify_unrestricted_good_reduction.py} は符号平方の階数2と行・列和、
二十点、180独立なSchur微分を確認する。さらに76個の係数に剰余二桁を指定し、
$53^3$ を法として独立な180式でNewton補正した結果が全200式を満たすことを確認する。
全ての精度と全ての次数に対する結論は上のエタール性と命題\ref{prop:unrestricted-good-reduction}による。

自由係数にその剰余類の有理整数を取れば、エタール射の零次元ファイバーなので
得る係数は $\mathbb Q$ 上代数的である。この整数格子は標数零の76次元空間に稠密である
（一変数ずつ多項式の零点を数えればよい）。局所環は相対滑らかで、その完備化は
$\mathbb Z_{53}[[z_1,\ldots,z_{76}]]$ である。剰余点を通る成分は一意で、
標数零の上記Walsh持上げを含む。この局所成分は幾何学的にも一意である。
実際係数体の任意の有限拡大に対しその付値環上で剰余点を固定すれば、
同じ完備局所環の記述がその付値環の係数で成立する。
よって格子の持上げの閉包は次元から定理\ref{thm:full-fifteen-component}の成分全体である。
同定理の一般の三条件は構成可能で稠密開集合を含み、その開への制限と四基底の商で主張を得る。
\end{proof}
